\section[B. AIMC Lab]{\myboxii{      \textbf{\texttt{B}}} AIMC Lab}

\textit{Time limit}   2 seconds

\textit{Memory limit} 1 Gigabyte

\textit{Input} \qquad standard input

\textit{Output} \qquad standard output


\begin{quote}
	\textit{街灯璀璨，可我日渐透明。}
\end{quote}

\subsection{Problem Statement}

Your AIMC Lab has developed an entirely new three-dimensional quantum matrix!

Consider a three-dimensional grid of size $N \times M \times H$, with coordinates $1 \le x \le N$, $1 \le y \le M$, and $1 \le z \le H$. A qubit occupies every integer grid point. Let $(x_u,y_u,z_u)$ be the coordinates of qubit $u$.

For any two qubits $u$ and $v$, define their \textbf{squared spatial distance} as
$$
d^2(u, v) = (x_u - x_v)^2 + (y_u - y_v)^2 + (z_u - z_v)^2.
$$

The system has eight energy matrices $A_{D,c}$, each of size $N \times M \times H$, indexed by integers $0 \le D \le 3$ and $0 \le c \le 1$.

You may initialize each qubit to either basis state $|0\rangle$ or $|1\rangle$. Its corresponding basis-state value is $c_P=0$ or $c_P=1$.

After initialization, the total energy is calculated as follows:

\begin{itemize}
	\item For each integer pair $0 \le D \le 3$, $0 \le c \le 1$, qubit $P$ gains $A_{D,c}(x_P,y_P,z_P)$ energy if and only if:
	\begin{itemize}
		\item The basis-state value of $P$ satisfies $c_P=c$.
		\item For every qubit $Q$ with $d^2(P,Q)=D$, the value $c_P+c_Q+D$ is even.
	\end{itemize}
\end{itemize}

As the chief scientist of AIMC Lab, initialize every qubit to a basis state so that the system's total energy is maximized. Output this maximum energy.

\subsection{Input}

The first line contains three integers $N$, $M$, and $H$ ($1 \le N,M,H \le 20$), the grid dimensions.

The data for the eight energy matrices $A_{D,c}$ follow in lexicographic order of $(D,c)$.

Each matrix consists of $N$ blocks. Each block has $M$ lines, each containing $H$ integers. Blocks are separated by blank lines.

The $z$-th number in the $y$-th line of the $x$-th block is the energy at $(x,y,z)$. All energy values satisfy $A_{D,c}(x_P,y_P,z_P) \in [0,10^9]$.

\subsection{Output}

Output one integer: the maximum total energy the system can achieve.

\subsection{Sample 1}

\begin{minipage}{\linewidth}
	\begin{minipage}[t]{0.49\linewidth}
		\begin{lstlisting}[title=\texttt{standard input}]
1 1 1
3

5

100

200

1000

2000

50

60
\end{lstlisting}
	\end{minipage}
	\begin{minipage}[t]{0.49\linewidth}
		\begin{lstlisting}[title=\texttt{standard output}]
2265
\end{lstlisting}
	\end{minipage}
\end{minipage}


\subsection{Sample 1 Explanation}

The grid contains only one qubit, at $(1,1,1)$.

Initializing it to $|1\rangle$ gives $A_{0,1}(1,1,1)+A_{1,1}(1,1,1)+A_{2,1}(1,1,1)+A_{3,1}(1,1,1)=5+200+2000+60=2265$, which is optimal.

\subsection{Sample 2}

\begin{minipage}{\linewidth}
	\begin{minipage}[t]{0.49\linewidth}
		\begin{lstlisting}[title=\texttt{standard input}]
2 2 2
1 1
1 1

1 1
1 1

0 0
0 0

0 0
0 0

0 0
0 0

0 0
0 0

0 0
0 0

0 0
0 0

10 10
10 10

10 10
10 10

0 0
0 0

0 0
0 0

0 0
0 0

0 0
0 0

0 0
0 0

0 0
0 0
\end{lstlisting}
	\end{minipage}
	\begin{minipage}[t]{0.49\linewidth}
		\begin{lstlisting}[title=\texttt{standard output}]
88
\end{lstlisting}
	\end{minipage}
\end{minipage}


\subsection{Sample 2 Explanation}

Initialize all $8$ qubits of the $2 \times 2 \times 2$ grid to $|0\rangle$.

\begin{itemize}
	\item Each qubit gains $1$ energy from $A_{0,0}$, for a total of $8$.
	\item All qubits at squared distance $2$ are also in state $|0\rangle$, so each qubit gains $10$ from $A_{2,0}$, for a total of $80$.
\end{itemize}

The total energy is $8+80=88$. This can be shown to be the maximum.
